\section*{Chapitre \ref{ch:ll:websocket-rest-tls-and-json-at-speed} --- WebSocket, REST, TLS et JSON à grande vitesse}

\begin{solution}{exo:ll:websocket-rest-tls-and-json-at-speed:1}
$2 + 2 + 4 + 300 = 308$ octets (en-tête, longueur sur 16 bits, clé de masquage, charge utile)~; $2 + 8 + 4 + 70\,000 = 70\,014$ octets~; une
trame émise par le serveur n'est pas masquée, donc $2 + 100 = 102$ octets.
\end{solution}

\begin{solution}{exo:ll:websocket-rest-tls-and-json-at-speed:2}
Deux allers-retours (TCP, puis TLS) et le calcul de la poignée de main~: $2 \times 0.5 + 2.06 \approx \qty{3.1}{\milli\second}$ avec une
poignée de main complète, environ \qty{2.6}{\milli\second} avec une \omterm{def:ll:websocket-rest-tls-and-json-at-speed:tls}{reprise de session}. Sur le portable mesuré, le calcul de la poignée de main, que
partagent les deux extrémités, coûte plus cher que les allers-retours.
\end{solution}

\begin{solution}{exo:ll:websocket-rest-tls-and-json-at-speed:3}
Les données émises sans aller-retour peuvent être rejouées~: la RFC~8446 accorde l'aller-retour économisé «~au prix de certaines propriétés de sécurité~». Un
ordre rejoué par un attaquant ou par un intermédiaire qui retransmet serait exécuté deux fois~; les ordres ne doivent circuler qu'après la
fin de la poignée de main.
\end{solution}

\begin{solution}{exo:ll:websocket-rest-tls-and-json-at-speed:4}
Chaque clé complétée forme exactement un bloc de 64 octets qui ne dépend que de la clé, si bien que l'état du hachage après ce bloc est le même pour
tous les messages. Sans précalcul~: le hachage interne traite le bloc ipad et le message de 119 octets complété sur deux blocs (trois
compressions), le hachage externe le bloc opad et le condensé interne de 32 octets complété sur un bloc (deux)~: cinq au total. Avec le précalcul~: trois.
\end{solution}

\begin{solution}{exo:ll:websocket-rest-tls-and-json-at-speed:5}
6\,412\,320\,000\,000 exactement. En virgule flottante binaire, 64123.2 n'est pas représentable~; multiplié par $10^8$, il tombe un cheveu
en dessous ou au-dessus de l'entier, si bien que la troncature peut donner 6\,412\,319\,999\,999, et des prix qui ne diffèrent que par le dernier chiffre peuvent
se comparer égaux ou inégaux par accident. Analyser le texte décimal en un entier évite la question.
\end{solution}

\begin{solution}{exo:ll:websocket-rest-tls-and-json-at-speed:6}
Répondre sans attendre au ping par un pong portant la même charge utile, conserver le premier fragment, et ne transmettre la mise à jour de profondeur seulement
lorsque son fragment final arrive~; les trames de contrôle peuvent s'intercaler entre les fragments, mais jamais les couper.
\end{solution}

\begin{solution}{exo:ll:websocket-rest-tls-and-json-at-speed:7}
Répéter la clé jusqu'à la longueur de la charge utile, puis convertir les deux chaînes d'octets en entiers
(\texttt{int.from\_bytes})~; effectuer le OU exclusif des deux entiers et reconvertir le résultat en octets (\texttt{to\_bytes}). Les tests sur fixtures ne changent pas~; le coût par trame de 300 octets passe de
quelques dizaines de microsecondes à quelques microsecondes.
\end{solution}

\begin{solution}{exo:ll:websocket-rest-tls-and-json-at-speed:8}
Chaque ordre paie deux allers-retours et une poignée de main, soit environ \qty{3.1}{\milli\second} ici~; la plateforme voit une nouvelle connexion par ordre
(et peut en limiter ou en bloquer l'usage)~; et le régulateur de limite de débit n'a aucun point unique où compter les requêtes. Une réserve de \omterm{def:ll:websocket-rest-tls-and-json-at-speed:rest}{connexions persistantes}
partagées derrière le régulateur évite les trois~: l'attente pour obtenir une connexion se compte en nanosecondes, la poignée de main en
millisecondes.
\end{solution}

\begin{solution}{pb:ll:websocket-rest-tls-and-json-at-speed:1}
\begin{enumerate}
\item Sur la mesure figurant dans le dépôt, environ \qty{0.9}{\micro\second} avec les composants C++ et \qty{48}{\micro\second} avec la
référence Python.
\item C++~: le décodage JSON (\qty{0.39}{\micro\second}). Python~: la trame octet par octet, \qty{16}{\micro\second} pour démasquer
et \qty{15}{\micro\second} pour masquer.
\item $0.93/(0.93 + 500) \approx 0.2~\%$ et $47.6/(47.6 + 500) \approx 8.7~\%$.
\item Environ \qty{16.9}{\micro\second}, soit $3.3~\%$ du chemin~: la mise en trame représentait les deux tiers du coût Python.
\item $2 \times \qty{0.5}{\milli\second}$ d'allers-retours plus la poignée de main~: environ \qty{3.1}{\milli\second} (complète) ou
\qty{2.6}{\milli\second} (\omterm{def:ll:websocket-rest-tls-and-json-at-speed:tls}{reprise de session}) par ordre.
\item Trois allers-retours et la poignée de main, soit environ \qty{3.6}{\milli\second}.
\item Ouvertes l'une après l'autre, environ 330 par seconde~; sur une \omterm{def:ll:websocket-rest-tls-and-json-at-speed:rest}{connexion persistante}, c'est la limite de débit de la plateforme qui s'applique.
\item Prévoir une reconnexion avant d'atteindre la limite~: ouvrir la nouvelle connexion, s'abonner et resynchroniser le carnet, puis fermer l'ancienne connexion
de sorte qu'aucune mise à jour ne soit manquée.
\item Il porte sur le prix et la quantité de l'ordre, connus seulement après la décision. L'état de hachage de la clé, les parties fixes de la
charge utile et la mise en forme de l'horodatage peuvent être préparés à l'avance.
\item 5\,000 mises à jour par seconde~: environ 0.2~\% d'un cœur avec le décodeur indexé et 4.9~\% avec la référence Python.
\item Mises à jour manquées~: un premier identifiant supérieur au dernier appliqué plus un. Le carnet est alors abandonné et reconstruit à partir d'un
nouvel instantané (\texttt{firm.wsbook} du livre~3).
\item Pour qu'une requête capturée ne puisse pas être rejouée plus tard~: la plateforme refuse celle qui est plus ancienne que sa fenêtre.
\item \textbf{Résultat nommé.} Avec un trajet de \qty{250}{\micro\second} dans chaque sens, le déchiffrement, la mise en trame, le JSON et la signature prennent 0.2~\% du
chemin plateforme--hôte--plateforme avec les composants C++ et 8.7~\% avec la référence Python~; réutiliser une connexion fait gagner environ
\qty{3.1}{\milli\second} par ordre REST par rapport à l'ouverture d'une connexion (\qty{2.6}{\milli\second} avec la \omterm{def:ll:websocket-rest-tls-and-json-at-speed:tls}{reprise de session}), soit plusieurs fois
le trajet réseau lui-même.
\item Lorsque les décisions de la stratégie valent des millisecondes et non des microsecondes, et que les messages sont peu nombreux~: le saut réseau et la
latence propre de la plateforme dominent de toute façon.
\item Dans un fil d'exécution ou un processus distinct qui termine le TLS pour plusieurs connexions, ou dans la prise en charge du TLS par le système d'exploitation ou par la
carte réseau lorsqu'elle existe, de sorte que le fil de la stratégie ne voie que des trames en clair.
\item Les mêmes coûts sur le processeur du serveur (ses instructions AES et SHA), pour les débits et les tailles de messages réels, ainsi que la
distribution du saut réseau, et non seulement sa médiane.
\item Ce n'est pas du temps processeur sur cet hôte, mais c'est la majeure partie du tick-to-trade et de sa variance~; un placement plus proche modifie
le résultat plus que tout décodeur.
\item Le JSON et le texte décimal, ainsi que le masquage~; l'analyse devient une lecture à décalages fixes (chapitre~16).
\item Alimenter le pool avec une séquence scénarisée de réponses de la plateforme comportant des 429 et un délai de nouvelle tentative, comme le régulateur du livre~3 les accepte par
l'intermédiaire de \texttt{on\_status}, et vérifier qu'il cesse d'envoyer jusqu'à l'écoulement du délai.
\item Il garde les connexions ouvertes, décode sans rien construire et signe sans recalculer le hachage de la clé.
\end{enumerate}
\end{solution}

\begin{solution}{iq:ll:websocket-rest-tls-and-json-at-speed:1}
Une nouvelle connexion coûte un aller-retour TCP, un aller-retour TLS et le calcul cryptographique de la poignée de main (des millisecondes), et peut être comptée
dans les limites de connexions de la plateforme~; une \omterm{def:ll:websocket-rest-tls-and-json-at-speed:rest}{connexion persistante} ne paie cela qu'une fois. Utiliser une réserve derrière le régulateur de limite de débit.
\iqlookfor{allers-retours plus calcul de la poignée de main, et limites.}
\end{solution}

\begin{solution}{iq:ll:websocket-rest-tls-and-json-at-speed:2}
Deux octets d'en-tête (bit final, code d'opération, bit de masquage, longueur sur 7 bits), une longueur étendue sur 16 ou 64 bits, une clé de masquage de 4 octets sur les trames du client,
puis la charge utile. Le masquage protège les mandataires intermédiaires contre des charges utiles fabriquées~; il coûte un OU exclusif par octet, soit des nanosecondes lorsqu'il est fait
par blocs de huit ou de trente-deux octets.
\iqlookfor{structure, raison, et coût bien réalisé.}
\end{solution}

\begin{solution}{iq:ll:websocket-rest-tls-and-json-at-speed:3}
Repérer tous les caractères structurants en une seule passe vectorisée, puis parcourir cet index selon la forme attendue du message, en lisant
les nombres sur place et en analysant les prix en entiers~; revenir à un analyseur général si la forme diffère. Pas d'arbre, pas d'allocation,
pas de virgule flottante.
\iqlookfor{un index structurel et un parcours propre à la forme.}
\end{solution}

\begin{solution}{iq:ll:websocket-rest-tls-and-json-at-speed:4}
HMAC-SHA-256 de la chaîne de requête et du corps avec la clé secrète, envoyé en hexadécimal avec la clé d'API, un horodatage et une fenêtre de validité
dans la charge utile. Hacher une seule fois les deux blocs complétés de la clé et recopier ces états à chaque requête~; utiliser une bibliothèque SHA-256 qui exploite les
instructions du processeur~; préparer les parties fixes de la charge utile.
\iqlookfor{le précalcul de la RFC 2104, et la protection contre le rejeu.}
\end{solution}

\begin{solution}{iq:ll:websocket-rest-tls-and-json-at-speed:5}
La \omterm{def:ll:websocket-rest-tls-and-json-at-speed:tls}{reprise de session} évite le certificat et une partie de l'échange de clés en réutilisant une clé d'une connexion antérieure~: moins de calcul,
mais toujours un aller-retour. Les données émises sans aller-retour partent avec le premier message et peuvent être rejouées, elles doivent donc être idempotentes,
ce qu'un ordre n'est pas.
\iqlookfor{calcul économisé face aux allers-retours, et rejeu.}
\end{solution}

\begin{solution}{iq:ll:websocket-rest-tls-and-json-at-speed:6}
La mise en place de la connexion sur le chemin de l'ordre (une nouvelle connexion ou une reconnexion par ordre), l'attente derrière un limiteur de débit ou derrière une
connexion occupée, le \omterm{def:ll:managed-and-functional-languages-on-the-desk:gc}{ramasse-miettes} ou un décodage lent côté client, un TLS ou une signature mal réalisés, et la latence propre de la plateforme~;
mesurer chaque étape avec des horodatages à chaque frontière.
\iqlookfor{mesurer les étapes~; les poignées de main d'abord.}
\end{solution}
